% CV body, shared by the résumé and the academic CV. Written with semantic
% macros only; style-sheets.tex decides how they look.
%
%   \ifacademic ... \fi       full academic CV only
%   \ifresume ... \fi         industry résumé only
%   \Todo{...}                missing information, printed in red
%
% Publications come from pubs.tex (generated from the website).
% Use real dashes (– —) rather than -- / ---.

\MakeHeader

\Summary{%
  Robot learning researcher and PhD candidate at the Intuitive Robots Lab, Karlsruhe Institute of
  Technology, advised by Rudolf Lioutikov. I give robot policies interpretable building blocks: concepts
  for what they perceive, primitives for how they move. Papers at CoRL, CVPR,
  ICRA, RSS and NeurIPS; reviewer for the major robot learning and machine learning venues.}

\Section{Experience}

\Entry{2022 – now}{PhD Researcher}{Intuitive Robots Lab, Karlsruhe Institute of Technology}{Karlsruhe, Germany}{%
  \begin{Highlights}
    \item Research on concept-aware robot policies: interpretable building blocks for what robots perceive
      (visual affordances; prototypes such as HyperPG, probabilistic prototypes on the hypersphere) and how
      they move (force-aware movement primitives; libraries of demonstrated motions, MPQ).
    \item One of the lab's first PhD researchers, helping build it up since its founding in 2022, including
      its bimanual robot learning stack: microservices from WebXR teleoperation and data processing to
      policy deployment on Franka Research 3 Duo.
    \item Papers at CoRL, CVPR, ICRA, RSS and NeurIPS, including datasets and benchmarks for robot vision
      and bimanual manipulation.
    \ifresume
    \item Teaching assistant for the Explainable AI lecture; co-teaches seminars and lab courses on
      explainable AI, behavior foundation models and robot learning; supervised 9 Bachelor's and Master's
      theses, one published at ICRA 2025.
    \fi
  \end{Highlights}}

\Entry{2025}{Intern, MLOps \& Robot VLA}{Poke \& Wiggle}{Munich, Germany}{%
  \begin{Highlights}
    \item Helped organize and host the Humanoid Manipulation \& AI Hackathon at the Deutsches Museum
      Munich, together with TUM.
    \item Implemented the hackathon's robot learning stack end to end: raw data processing, policy
      training and deployment.
  \end{Highlights}}

\Entry{2013 – 2017}{Research Working Student}{SAP, Big Data Intelligence and Research}{}{}

\Section{Education}

\Entry{2018 – 2022}{M.Sc. Computer Science}{Karlsruhe Institute of Technology}{Karlsruhe, Germany}{%
  \begin{Highlights}
    \item Thesis: \emph{Imitation Learning for Versatile Robot Manipulation} (Autonomous Learning Robots Lab).
      Research assistant at the Human-Centered Systems Lab and the Autonomous Learning Robots Lab.
  \end{Highlights}}

\Entry{2013 – 2018}{B.Sc. Computer Science}{TU Darmstadt}{Darmstadt, Germany}{%
  \begin{Highlights}
    \item Thesis: \emph{Development of a Content-Based Image Retrieval System via Deep Learning}
      (Fraunhofer SIT). Exchange student at Tongji University, Shanghai.
  \end{Highlights}}

\ifresume
  \Section{Selected Publications}
  \PubList{selected}{}
  \MorePubs
\fi
\ifacademic
  \Section{Publications}
  \PubList{conference}{Conference papers}
  \PubList{journal}{Journal articles}
  \PubList{workshop}{Workshop papers}
  \PubList{preprint}{Preprints}
\fi

\Section{Talks \& Service}

\ifresume
\SkillRow{Talk}{\emph{Towards Concept-Aware Robot Policies}, Shanghai Research Institute for Intelligent
  Autonomous Systems, Tongji University, Apr 2026}
\else
\Entry{Apr 2026}{Towards Concept-Aware Robot Policies}{Shanghai Research Institute for Intelligent Autonomous
  Systems, Tongji University}{Shanghai, China}{}
\fi
\SkillRow{Reviewer}{CoRL, RSS, ICRA, NeurIPS, ICML, ICLR, CVPR, WACV, JMLR}

\ifacademic
\Section{Teaching}

\Course{Explainable Artificial Intelligence}{Lecture \textperiodcentered{} Teaching assistant}{%
  SS 2023 – SS 2026. Concepts and methods of explainable AI from the perspective of robot learning and
  human–robot interaction.}
\Course{Explainable Artificial Intelligence}{Seminar \& research project}{%
  WS 23/24 – WS 25/26. Students study XAI methods for robot learning, then implement and evaluate them.}
\Course{Behavior Foundation Models}{Seminar \& research practical course}{%
  SS 2026. Students study and develop methods for robot learning with behavior foundation models, such
  as VLAs.}
\Course{Robot Learning}{Bachelor's lab course \& proseminar}{%
  2024 – 2026. Introduction to PyTorch, MuJoCo benchmarks and the tools of robot learning research.}

\Section{Thesis Supervision}
\Note{Supervised 9 Bachelor's and Master's theses, including:}

\Thesis{M.Sc.}{Force-Aware Robot Manipulation using Movement Primitives}{Paul Lödige}{published at ICRA 2025}
\Thesis{M.Sc.}{Concept Motivated Prototypes: Learning Concept Representations with Probabilistic Prototypes}{Korbinian Rudolf}{}
\fi

\ifresume\Section{Skills \& Interests}\else\Section{Skills}\fi

\SkillRow{Robot learning}{Python, PyTorch, LeRobot, starVLA, MuJoCo, Weights \& Biases}
\SkillRow{Robotics \& infrastructure}{ROS~2, Docker, Franka Panda, Franka Research 3 Duo, WebXR}
\SkillRow{Languages}{German (native), Chinese (bilingual), English (professional)}

\ifresume
\SkillRow{Interests}{Anything with a good story: books, video games, and running tabletop campaigns as a
  Dungeons \& Dragons game master.}
\else
\Section{Outside the Lab}
\Plain{Anything with a good story: books, video games, and running tabletop campaigns as a
  Dungeons \& Dragons game master.}
\fi
